\section*{Заключение}

Предложенный подход к выявлению вредоносной активности на основе атрибутивных метаграфов демонстрирует высокую эффективность за счёт:
\begin{enumerate}
    \item совместного учёта структурных и семантических характеристик событий;
    \item естественной интеграции с MITRE ATT\&CK через паттерны;
    \item возможности контекстной фильтрации, снижающей ложные срабатывания.
\end{enumerate}
Метод может быть реализован как модуль для существующих SIEM/SOAR-платформ. В частности, он позволяет:
\begin{enumerate}
    \item автоматически сопоставлять события с тактиками MITRE ATT\&CK;
    \item строить цепочки рассуждений (kill chain) при расследовании инцидентов;
    \item генерировать рекомендации по реагированию на основе контекста атаки.
\end{enumerate}
В дальнейшем планируется:
\begin{enumerate}
    \item расширение библиотеки паттернов за счёт MITRE D3FEND (методов защиты);
    \item поддержка онтологий, таких как STIX 2.1, для совместимости с платформами анализа угроз.
\end{enumerate}
